\chapter{Conclusions \& Future Work}
\label{ch:conclusions}

\section{Conclusions}
\label{sec:conclusions}

This thesis asked whether a controller can adapt to the body it is given at a cost small enough to sit inside a search over bodies, and whether such an adaptation changes what the search finds. The shape of any answer was fixed in \autoref{ch:challenges}: a controller trained once for the whole space of bodies, and an adaptation to the body being scored that costs no more than the flights that score it. That shape was given a content in \autoref{ch:method}: a generalist trained by reinforcement learning and then frozen, Hebbian rules on its output layer evolved by CMA-ES, and a pipeline that nests the search of the rules inside a multi-objective search over the bodies. In \autoref{ch:results} the pipeline was run eight times against eight runs of the search under the frozen generalist, and the inner loop was isolated on fixed bodies. \emph{The experiments answer the question in the affirmative: an adaptation that costs no more than the flights that score a body changes the outcome of the search over bodies, and it changes that outcome for the better on both objectives.} The paragraphs that follow state what was found, in the order in which the thesis built it.

\paragraph{The adaptation costs no more than the flights.} The whole pipeline consists of flights of a network that is never trained again. A run of 200 generations flies about a quarter of a million flights per generation, examines 3\,136 bodies, of which at most 1\,600 are distinct, and takes about four days on one computing node with two graphics cards. The control condition, which updates its bodies at every generation and gives each body the same number of exam flights, took about 130 hours against 94. A controller training for each distinct body would have taken almost two months with a specialist trained for each body, and two and a half years with a controller of the generalist's kind. The cost of the adaptation is the cost of the flights, as the computational feasibility challenge required.

\paragraph{Co-design improves on morphology-only evolution on both objectives.} On the exam course the cumulative fronts of the two conditions agree at their two ends and separate in between. The tip of a front is common to both conditions: no body of either condition flies past 226~m of the 300~m course. The cheap arm of the front, the bodies that fly 85 to 110~m for a cost of transport of 0.10 to 0.12, is the same in both conditions. Between 128 and 213~m of progress the co-design fronts are cheaper at every level. At 150, 170 and 190~m the cheapest co-design body spends 16 to 17\,\% less energy per metre than the cheapest body of the control. Within a cost of transport of 0.15 to 0.25 the farthest co-design body flies 9 to 12\,\% farther. On each of these indicators every one of the eight co-design runs is better than every one of the eight control runs, and the hypervolume of the fronts separates the two groups as completely. The control was given more exams, more examined bodies and more computing time, and it did not catch up. The reference drone, a body of ordinary proportions, lies on the mean front of the control. At the progress of the reference drone the control finds bodies that are 4 to 6\,\% cheaper, and the co-design runs find bodies that are 18 to 23\,\% cheaper. \emph{The adaptation makes the search find bodies that spend less at equal progress and travel farther at equal cost of transport, and the margin over the reference drone is several times the margin that the search under the fixed controller obtains from the same starting bodies.}

\paragraph{The body search takes a different direction when the controller adapts.} Started from the same 64 bodies, the two conditions do not go to the same place in the space of the genome. For three of the six seeds that carry a run of each kind, the two runs end farther apart than control runs do among themselves. For the other three seeds they end as far apart as two control runs started from different bodies. Seven of the eight morphology-only runs end in one cluster, and the eighth ends elsewhere. Under a fixed controller the search over bodies has one main optimum, and most runs find it. The co-design runs end farther apart from each other, and farther from the cluster than the clustered control runs are from one another. A co-design run is no closer to the control run that started from its own bodies than to a control run that started from other bodies. This is the morphology optimization challenge observed in the data. \emph{An adapting controller does not shift the optimum of the body search by a fixed amount: it changes which bodies come first.} The two ends of one co-design front show how different the bodies that come first can be. One end is a light body of 0.58~kg with wing panels of 0.40~m and its centre of mass well aft, which flies 220~m at a cost of transport of 0.39. The other end is a glider of 1.96~kg with a wing of almost a square metre and its centre of mass forward, which flies 128~m at a cost of transport of 0.13. The two bodies share the gear ratios of their joints and a cambered airfoil: the run moved the wing and the fuselage, and kept the actuation.

\paragraph{The controller adapts to the morphology.} A co-design front is a front of pairs. Flown again under their own rules, the bodies of a co-design front travel 9 to 19~m farther than under the frozen generalist, and almost every body of every front travels farther. The same rules transplanted onto the bodies of a morphology-only front add 1 to 6~m. The progress that the rules add therefore belongs to the bodies they were evolved with, and those bodies were retained because the rules fly them. A co-design front is not the front of the generalist with a better controller on top. The weights that the rules produce depend on the body. The same set of rules, flown on the two bodies at the ends of a front, arrives at weight changes that differ most on the two sweep commands. The sweep is the command by which the two bodies differ most in what they can do. The rules are never told which body they fly. The body reaches them through the activity it induces in the frozen layers, and the rules turn that activity into gains that suit that body. \emph{This is adaptation to the body in the sense in which \autoref{sec:hebbian-controller} defined it: the same rules, given a different body, arrive at a different controller.}

\paragraph{The rules recover part of what the generalist lacks on a given body.} On the reference drone the frozen generalist scores 7.7 points of reward less than a specialist trained for that drone, a gap of 4\,\%. Rules evolved on that body recover 3.4 points, 44\,\% of the gap; measured against the three specialists that were trained, the recovery lies between a quarter and a half of the gap. The two gains are of a different kind. The advantage of the specialist lies in avoiding the trees. The specialist loses 28\,\% of its flights where the generalist loses 53\,\%, and it travels farther at a higher cost of transport. The rules leave the crashes and the progress where the generalist has them, and they lower the cost of transport by 9\,\% at a speed 3\,\% lower. In reward, the rules recover part of what the generalist lacks. In objective space, they move the body along the cost of transport at constant progress, an axis along which the specialist does not move. The same decomposition holds for rules evolved on 64 random bodies and measured on the reference drone, which that search never flew: a smaller gain, of the same kind.

\paragraph{On a specialist, the rules make the flight dynamics plastic.} Evolved on top of a specialist, the rules add more than anywhere else on a fixed body: 9.6 and 9.9 points on two of the three specialists, over 5\,\%, and nothing on the second. The frozen controllers explain the difference. The three specialists were trained with the same settings and different seeds, and they did not arrive at the same flight of the reference drone. The first and the third fly it at 13.9~m/s with a high cost of transport; the second flies it at 13.3~m/s and cheaply. Under their rules, all four base controllers, the generalist included, end at the same flight: 13.3 to 13.4~m/s, with crashes and progress untouched and a lower cost of transport. The second specialist was flying that flight already, and the rules leave it there. \emph{Where a base controller ends up under the rules does not depend on where it started, and what it gains is the distance it had to travel.} A specialist with rules is, in effect, a controller whose flight dynamics are plastic. The gains of its output layer, and with them the dynamics of the closed loop that the body forms with its controller, are retuned during the flight towards the flight that the fitness rewards on this body. The gain is real in the currency of the search, because a lower cost of transport at equal progress is a better point in objective space.

\paragraph{The price is the generality of the controller.} Along a run, the best rules of each generation fly the bodies of the final front better and better, up to 17\,\% more fitness than the frozen generalist, and they fly bodies drawn at random from the design space worse and worse, down to 18\,\% less. The reference drone, which is never in the population, ends 6 to 9\,\% worse under the rules than under the frozen generalist. The rules specialize to the family of bodies that the outer loop retains, and they could not be reused on another family without being evolved again. For the search this is not a defect. The search asked for an adaptation to the bodies being scored, and the adaptation to a family rather than to every body in the space is what keeps the adaptation cheap. For any other use of the evolved rules it is a limit.

\paragraph{Limits of the study.} Every quantity of this thesis was measured in simulation, with an aerodynamic model that is trusted for the order of the bodies rather than for absolute values. The tip of the fronts is the same in both conditions, so the comparison was made where the bodies had room to differ. The fixed-body gains were measured on one reference drone and on three specialists whose training reached different flights with the same settings, so that what the rules add to a controller depends on where that controller starts. One reward function was used as the fitness of the rules, and the economy that the rules find is the change that this reward pays for. The experiments say what an adaptation this cheap does inside this search. They do not say that evolved Hebbian rules are the best cheap adaptation, nor that the fronts found are the best fronts of the design space.

\emph{Within these limits, the thesis shows that a controller trained once and a few hundred evolved coefficients are enough to let a search over bodies find bodies that the adapted controller flies better than the fixed one. The results support the hypothesis that a search under the fixed controller would have discarded those bodies.}

\section{Future Work}
\label{sec:future-work}

Three directions follow from the results. The first concerns the representation on which the rules act, which this thesis left as the training produced it. The second and the third concern the transfer of the mechanism, to real drones and to systems other than drones.

\subsection{Meta-Learning Techniques}
\label{subsec:future-meta-learning}

The rules act on the latent representation of the last hidden layer. The frozen layers compute 32 features from the history of observations and commands, and the rules change only the gains that translate those features into the commands. The quality of the adaptation is therefore bounded by the quality of the representation. In this thesis the representation was trained by reinforcement learning for a controller with fixed gains, and nothing in that training anticipated that the gains would become plastic. The features are those that serve one compromise set of gains for all bodies, not those that would serve a set of gains adapted to each body. Gradient-based meta-learning trains representations for exactly this situation. MAML \cite{finn2017maml} learns an initialization from which a few gradient steps adapt a network to a new task. Raghu et al.\ \cite{raghu2020rapid} showed that adapting the last layer alone is enough when the features below are trained with that adaptation in the loop. If the bodies are read as tasks, the same procedure would train the features of the generalist so that an adaptation of the output layer to a body pays as much as possible. The Hebbian rules would then act on a representation built to be read by an adapted output layer. Two forms of this are possible. In the first, the features are meta-trained with gradient steps on the output layer over the training bodies, and the rules are evolved afterwards, as they are now. In the second, the rules of an earlier run are kept in place during the training of the generalist, so that the features are trained with the plastic output layer in the loop. The training of the features and the evolution of the rules then alternate. The second form resembles the co-training of differentiable plasticity \cite{miconi2018differentiable}, with the rules evolved instead of differentiated. \emph{The expected gain is in the half of the mechanism that this thesis kept fixed: the representation that the rules read.}

\subsection{Testing on Real Drones}
\label{subsec:future-real-drones}

The rules adapt the control to the body that is flying, read through the answer of that body to the commands. A real drone differs from its simulated model in its mass and inertia, in the aerodynamic forces on its surfaces and in the response of its servomotors. From the point of view of the rules, these differences are another body. The adaptation that the rules provide to a body they never met may therefore extend to the gap between simulation and reality, during flight and without a model of the discrepancy. The literature on evolved plastic networks gives reasons to expect this. Plastic controllers evolved in simulation transferred to a physical robot better than controllers with evolved fixed weights \cite{floreano2000evolutionary}. Shared Hebbian rules transferred to ten physical robots with a loss of under one percent \cite{vandiggelen2025emergent}. The base controller of this thesis was trained with the randomization of forces, masses and latencies on which transfer from simulation usually relies \cite{tobin2017domain, peng2018sim}. The rules would add an online adaptation to that offline robustness. None of this can be shown in simulation. The test requires that a body of a front, or the reference drone, be built with its morphing joints, that the network and the Hebbian update run on the onboard computer, and that the depth sectors be obtained from an onboard sensor. The computational part is within reach of any flight controller: the update of the 224 weights is a few thousand arithmetic operations per control step, at 25 steps per second. The laboratory where this work was carried out flies morphing winged drones of the same family \cite{bergonti2024codesign}, so the platform exists. \emph{Whether the rules close part of the gap between simulation and reality, leave it unchanged or widen it is a question that only a flight can answer.}

\subsection{Testing on Different Systems}
\label{subsec:future-other-systems}

Nothing in the mechanism is specific to drones. It requires a network pretrained for a task over a family of conditions, a linear output layer, a stream of activity through that layer, and a quantity that varies between lifetimes and that the network is not told. Hebbian rules on top of a pretrained controller are, in this sense, a system-agnostic model. Within robotics, the body is not the only quantity of this kind. A payload, a worn actuator or a damaged limb change the dynamics in the same way, and the robustness to damage of evolved plastic networks is the result that made them known \cite{najarro2020meta}. Outside robotics, the same structure appears wherever a pretrained model follows a stream whose generating process changes between episodes. Sequence analysis is one case: a model trained on many sources is applied to a new source, with no labels on which to fine-tune it. The analysis of stochastic processes is another: the parameters of the process are fixed within an episode, hidden from the model and different from one episode to the next. This is the situation that this thesis described for the body. Differentiable plasticity has been applied to memorization and few-shot classification \cite{miconi2018differentiable}, so plastic weights can act on tasks of that kind. The open question is whether evolved rules on top of a frozen pretrained model give there the gain they gave here. The experiments of this thesis also say what to measure. \emph{The gain found here was a specialization to the family that the search retained, not a general improvement, so the measurement in any new system must be made both on the family and outside it.}
